\editionPage[\pdfbookmark[0]{Deux brevets blancs}{bm-bilans-guide-parent}]{bilans-guide}{Deux brevets blancs}{Session 2029 · Sujets originaux}{Durée totale : 2 heures · Note sur 20 · Coefficient 2}
\begin{blocEdition}{bilans-guide}

Définition retenue : note de service du 11 septembre 2026 (BO spécial n° 4 du
17 septembre 2026), en vigueur à compter de la session 2027. Au 23 octobre 2026,
aucun texte propre à la session 2029 n’était publié ; vérifier avant l’examen
qu’aucune modification n’est intervenue.
\editionEtape{Partie 1 : 20 minutes, sans calculatrice}
Automatismes et applications directes : 6 points. Questions à réponses courtes,
portant sur les automatismes du cycle 4 ou l’application directe d’une propriété,
d’une formule ou d’un théorème. Ici, douze questions à 0,5 point. La première
copie est relevée au bout de 20 minutes.
\editionEtape{Partie 2 : 1 h 40, calculatrice autorisée}
Raisonnement et résolution de problèmes : 14 points. Exercices indépendants,
sans QCM ; toutes les réponses sont justifiées, sauf indication contraire. La
maîtrise de la langue est évaluée sur 2 points. Utiliser une seconde copie ;
le passage de la première à la seconde partie se fait sans pause.
\editionEtape{Le barème de ce livre}
Chaque problème vaut 3,5 points : trois questions à 1 point et 0,5 point pour
la maîtrise de la langue et la clarté de la rédaction. Les 2 points de maîtrise
de la langue sont ainsi inclus dans les 14 points. Ce partage est un choix
d’entraînement ; il ne préjuge pas des barèmes des sujets officiels.
\editionEtape{Après la séance}
Comparer les démarches, classer les erreurs par notion et refaire les questions
utiles sans correction. Les essais et démarches engagés, même non aboutis, sont
valorisés : un résultat numérique faux n’efface pas une démarche correcte.

\end{blocEdition}
\editionPage{bilan-A-auto}{Brevet blanc A}{Relier calculs, modèles et figures}{Partie 1 · 20 min · Sans calculatrice · 6 points}
\begin{blocEdition}{bilan-A-auto}
\begin{questions}
\item Calculer $-7+3\times5$.
\item Calculer $\dfrac23+\dfrac16$.
\item Donner la notation scientifique de 450000.
\item Résoudre $x^2=81$.
\item Développer $(x+3)^2$.
\item Résoudre l’inéquation $-2x\ge8$.
\item Calculer 15 \% de 80.
\item Pour $f(x)=2x+1$, calculer $f(-2)$.
\item Partager 35 € selon le ratio $2:5$.
\item Donner la médiane de $2\,;4\,;5\,;8\,;11$.
\item $P(A)=0{,}5$, $P(B)=0{,}4$ et $P(A\cap B)=0{,}2$. Calculer $P(A\cup B)$.
\item ABCD est un parallélogramme. Compléter : $\vect{AB}=\vect{D\ldots}$.\end{questions}
\end{blocEdition}
\editionPage{bilan-A-ex-1}{Brevet A · Problème 1}{Un atelier de carnets}{Partie 2 · 1 h 40 pour les quatre problèmes · 3,5 points}
\begin{blocEdition}{bilan-A-ex-1}
Un atelier possède 168 feuilles bleues et 126 feuilles blanches. Il prépare le plus grand nombre possible de carnets identiques, sans reste. Chaque carnet contient le même nombre de feuilles de chaque couleur.\begin{questions}
\item Déterminer le nombre maximal de carnets et leur composition.
\item L’atelier vend chaque carnet 4,50 €. Les fournitures ont coûté 78 €. Calculer le solde si tous les carnets sont vendus.
\item Un autre atelier vend le carnet 5 €, puis applique 10 \% de réduction. Comparer son prix unitaire au précédent.
\end{questions}\par\medskip{\small\color{otMuted}Barème : question 1 : 1 pt ; question 2 : 1 pt ; question 3 : 1 pt ; maîtrise de la langue et clarté de la rédaction : 0,5 pt.}
\end{blocEdition}
\editionPage{bilan-A-ex-2}{Brevet A · Problème 2}{Une voile divisée}{Partie 2 · 1 h 40 pour les quatre problèmes · 3,5 points}
\begin{blocEdition}{bilan-A-ex-2}
ABC est rectangle en A, AB = 12 m, AC = 9 m. M appartient à [AB] avec AM = 4 m. N appartient à [AC] et (MN) est parallèle à (BC).\begin{questions}
\item Calculer BC.
\item Calculer AN et MN.
\item On construit le point D, image de N par la translation de vecteur $\vect{MB}$. Quelle est la nature du quadrilatère MBDN ? Démontrer que D appartient à la droite (BC).
\end{questions}\par\medskip{\small\color{otMuted}Barème : question 1 : 1 pt ; question 2 : 1 pt ; question 3 : 1 pt ; maîtrise de la langue et clarté de la rédaction : 0,5 pt.}\label{fig-bilan-A-ex-2}\figureGeo{brevet-A-geo}

\end{blocEdition}
\editionPage{bilan-A-ex-3}{Brevet A · Problème 3}{Deux locations de matériel}{Partie 2 · 1 h 40 pour les quatre problèmes · 3,5 points}
\begin{blocEdition}{bilan-A-ex-3}
Pour $n$ journées entières de location, A facture $18+7n$ euros et B facture $10n$ euros, pour $1\le n\le15$.\begin{questions}
\item Calculer les deux prix pour 4 jours. Quelle offre choisir ?
\item Pour combien de jours les prix sont-ils égaux ? Résoudre l’inéquation $18+7n<10n$ et interpréter le résultat.
\item Avec un budget de 90 €, déterminer la durée maximale possible chez chaque loueur, en résolvant une inéquation.
\end{questions}\par\medskip{\small\color{otMuted}Barème : question 1 : 1 pt ; question 2 : 1 pt ; question 3 : 1 pt ; maîtrise de la langue et clarté de la rédaction : 0,5 pt.}\label{fig-bilan-A-ex-3}\figureGeo{brevet-A-tarifs}

\end{blocEdition}
\editionPage{bilan-A-ex-4}{Brevet A · Problème 4}{Des scores et un tirage}{Partie 2 · 1 h 40 pour les quatre problèmes · 3,5 points}
\begin{blocEdition}{bilan-A-ex-4}
Une série de 20 essais donne les scores du tableau. On tire ensuite deux boules avec remise dans un sac de 3 boules rouges et 1 bleue, numérotées $R_1$, $R_2$, $R_3$ et $B$ ; les tirages sont faits au hasard.\begin{center}\renewcommand{\arraystretch}{1.35}\begin{tabular}{c|cccc}\toprule
Score & $0$ & $1$ & $2$ & $3$\\\midrule
Effectif & $2$ & $5$ & $8$ & $5$\\
\bottomrule\end{tabular}\end{center}\begin{questions}
\item Calculer la moyenne, la médiane et les quartiles des scores.
\item À l’aide du tableau des couples, calculer la probabilité d’obtenir deux boules rouges.
\item Calculer la probabilité d’obtenir exactement une boule bleue, puis celle d’obtenir au moins une boule rouge.
\end{questions}\par\medskip{\small\color{otMuted}Barème : question 1 : 1 pt ; question 2 : 1 pt ; question 3 : 1 pt ; maîtrise de la langue et clarté de la rédaction : 0,5 pt.}\label{fig-bilan-A-ex-4}\figureGeo{brevet-A-tableau}

\end{blocEdition}
\editionPage{bilan-B-auto}{Brevet blanc B}{Choisir une méthode et contrôler}{Partie 1 · 20 min · Sans calculatrice · 6 points}
\begin{blocEdition}{bilan-B-auto}
\begin{questions}
\item Calculer $18-4\times3$.
\item Calculer $\dfrac35\div\dfrac9{10}$.
\item Écrire $10^3\times10^{-5}$ comme une puissance de 10.
\item Décomposer 60 en produit de facteurs premiers.
\item Factoriser $x^2-49$.
\item Résoudre $(x-3)(2x+1)=0$.
\item Augmenter 50 de 12 \%.
\item Pour $g(x)=-3x$, donner l’antécédent de 12.
\item Donner le volume exact d’une boule de rayon 3 cm.
\item Donner le premier quartile de $1\,;3\,;4\,;6\,;8\,;9\,;12\,;15$.
\item Un évènement a une probabilité de 0,3. Donner celle de son contraire.
\item Simplifier $\vect{AB}+\vect{BC}$.\end{questions}
\end{blocEdition}
\editionPage{bilan-B-ex-1}{Brevet B · Problème 1}{Un abri et sa pente}{Partie 2 · 1 h 40 pour les quatre problèmes · 3,5 points}
\begin{blocEdition}{bilan-B-ex-1}
Un toit est représenté par ABC rectangle en A. AB = 4 m horizontalement, AC = 2,5 m verticalement et BC est le versant. L’abri a une longueur de 6 m perpendiculaire au dessin.\begin{questions}
\item Calculer BC au centimètre.
\item Calculer l’angle $\widehat B$ au degré.
\item Le versant rectangulaire doit être couvert. Ajouter 8 \% à son aire pour les découpes ; chaque plaque couvre 2 m$^2$. Combien de plaques au minimum ? Utiliser la longueur BC non arrondie.
\end{questions}\par\medskip{\small\color{otMuted}Barème : question 1 : 1 pt ; question 2 : 1 pt ; question 3 : 1 pt ; maîtrise de la langue et clarté de la rédaction : 0,5 pt.}\label{fig-bilan-B-ex-1}\figureGeo{brevet-B-geo}

\end{blocEdition}
\editionPage{bilan-B-ex-2}{Brevet B · Problème 2}{Une bouteille et une boule}{Partie 2 · 1 h 40 pour les quatre problèmes · 3,5 points}
\begin{blocEdition}{bilan-B-ex-2}
Une bouteille est modélisée par un cylindre de rayon intérieur 4 cm et de hauteur utile 20 cm. Une machine la remplit à 50 mL/s, à débit constant.\begin{questions}
\item Calculer la capacité exacte en cm$^3$, puis l’arrondir au mL.
\item Calculer la durée nécessaire pour remplir 90 \% de la capacité, au dixième de seconde.
\item On verse le contenu d’une bouteille pleine dans un récipient en forme de boule de rayon intérieur 7 cm. Le récipient est-il assez grand ? Quel pourcentage de son volume est rempli, à l’unité près ?
\end{questions}\par\medskip{\small\color{otMuted}Barème : question 1 : 1 pt ; question 2 : 1 pt ; question 3 : 1 pt ; maîtrise de la langue et clarté de la rédaction : 0,5 pt.}\label{fig-bilan-B-ex-2}\figureGeo{brevet-B-cylindre}

\end{blocEdition}
\editionPage{bilan-B-ex-3}{Brevet B · Problème 3}{Un script de figures}{Partie 2 · 1 h 40 pour les quatre problèmes · 3,5 points}
\begin{blocEdition}{bilan-B-ex-3}
Le projet démarre stylo baissé, orienté vers la droite. Le bloc carré(c) trace un carré de côté c et revient au départ avec l’orientation initiale. Le programme principal est donné ci-dessous.\begin{questions}
\item Donner les côtés successifs des carrés tracés et le nombre de passages dans la boucle.
\item Calculer la somme des périmètres tracés.
\item Donner la valeur affichée à la fin et la distance totale parcourue stylo relevé. Expliquer pourquoi cette valeur finale de c n’est pas le côté d’un carré tracé.
\end{questions}\par\medskip{\small\color{otMuted}Barème : question 1 : 1 pt ; question 2 : 1 pt ; question 3 : 1 pt ; maîtrise de la langue et clarté de la rédaction : 0,5 pt.}\label{fig-bilan-B-ex-3}\figureGeo{brevet-B-script}

\end{blocEdition}
\editionPage{bilan-B-ex-4}{Brevet B · Problème 4}{Un jeu à deux étapes}{Partie 2 · 1 h 40 pour les quatre problèmes · 3,5 points}
\begin{blocEdition}{bilan-B-ex-4}
Une roue possède quatre secteurs égaux : trois portent la lettre A, un porte la lettre B. Un sac contient une boule rouge et une bleue. On tourne la roue, puis on tire une boule au hasard, indépendamment de la roue. Le joueur gagne s’il obtient A et rouge, ou B et bleue. Les secteurs A sont numérotés A1, A2, A3 dans le tableau.\begin{questions}
\item Combien d’issues équiprobables compte l’expérience ? Calculer la probabilité de gagner.
\item Sur 100 parties, peut-on garantir exactement 50 gains ? Expliquer.
\item On note E : « la roue donne A » et F : « la boule est rouge ». Calculer $P(E)$, $P(F)$, $P(E\cap F)$, puis $P(E\cup F)$.
\end{questions}\par\medskip{\small\color{otMuted}Barème : question 1 : 1 pt ; question 2 : 1 pt ; question 3 : 1 pt ; maîtrise de la langue et clarté de la rédaction : 0,5 pt.}\label{fig-bilan-B-ex-4}\figureGeo{brevet-B-tableau}

\end{blocEdition}
